\section{Síntesis y ampliación}

Esta clase parte del despliegue heterogéneo y avanza capa por capa por la delivery de software, el runtime substrate, la decision chain, el workflow de manufactura y la retroalimentación institucional. Apollo somete el estado del software y del entorno a un control continuo; Rubix integra la seguridad y el ciclo de vida en un substrate unificado; el caso de Maven muestra que el modelo debe entrar en un OODA loop con permisos y retroalimentación; el caso de manufactura muestra que el workflow debe recalcularse cuando cambian las restricciones; y la privacy-security frontier exige que el aumento de capacidades y el debido proceso existan al mismo tiempo.

\begin{importantbox}{Nueve conclusiones transferibles}
\begin{enumerate}
  \item Al describir la scale, cuantifica al mismo tiempo la heterogeneidad de los entornos, la frecuencia de cambios y el tiempo sin conexión, no solo el número de nodos.
  \item Escribe el despliegue de software como un desired-state control loop, no como una serie de shell scripts no observables.
  \item Que el service owner declare health, dependency y schema compatibility; solo así la automatización puede ser segura.
  \item Usa release channel y canary para expresar el grado de confianza, en lugar de buscar que todos los entornos cambien de forma sincronizada.
  \item Compliance-as-code debe impedir estados que incumplen las normas y generar evidencia, no solo llenar automáticamente los formularios de auditoría.
  \item La Ephemeral infrastructure acorta la ventana de drift y de persistencia, pero exige que la aplicación admita failover.
  \item La inference del modelo es solo un eslabón de la decision chain; los permisos, el workflow y el outcome son los que determinan el valor del despliegue.
  \item Las spreadsheets en campo son evidencia de requisitos de producto y de cortes en la retroalimentación; no deben simplemente prohibirse.
  \item La tecnología puede ampliar el conjunto factible de privacy/security, pero los objetivos legítimos siguen requiriendo decisiones institucionales y supervisión.
\end{enumerate}
\end{importantbox}

\subsection{Ejercicio práctico: diseñar una actualización de seguridad para una Fleet heterogénea}

Supón que mantienes un producto de 30 servicios que abarca nube pública, centros de datos de clientes y sitios de borde sin conexión a la red, y que necesitas corregir una vulnerabilidad de alta gravedad en una library. Diseña la catalog metadata, los release channels, los health monitors, el rollback y el proceso de excepciones. El ejercicio no puede limitarse a escribir «usar Kubernetes» y «despliegue automático»; debe especificar las environment constraints y la compatibilidad de los datos.

\begin{enumerate}
  \item Enumera al menos cuatro tipos de entorno y define su conectividad, ventana de mantenimiento, nivel de seguridad y capacidad disponible.
  \item Para cada component afectado, crea registros de dependency, schema, vulnerability y owner.
  \item Diseña un rollout canary→connected→regulated→air-gapped y escribe las condiciones de entrada y salida de cada etapa.
  \item Elige un entorno que no pueda actualizarse de inmediato y especifica el compensating control, el owner del riesgo y la fecha de vencimiento.
  \item Simula una falla de la nueva versión y explica cómo determinar si el rollback es seguro y cómo manejar los datos que ya se escribieron.
\end{enumerate}

\begin{knowledgebox}{Criterios de aceptación}
Una propuesta de alta calidad debe poder responder «qué ejecuta ahora cada entorno», «por qué puede actualizarse», «a quién se notifica cuando algo falla», «quién aprueba las excepciones» y «qué evidencia demuestra que el riesgo disminuyó». Si las respuestas dependen de que un ingeniero con experiencia recuerde todo el orden de los pasos, el plano de control todavía no se ha formado realmente.
\end{knowledgebox}

\subsection{Lecturas adicionales}

Se recomienda leer directamente `source-materials/apollo-demo-day-transcript.pdf` en el directorio de esta clase y contrastarlo con los ejemplos de Apollo Catalog, release channel, monitor y vulnerability; después, leer la documentación oficial vigente de Palantir Rubix y comparar las 40--72 horas que se mencionaron oralmente en clase con la formulación actual de 48 horas de node-lifetime. Por último, elige un workflow de alto riesgo fuera del ámbito de la defensa, por ejemplo, las camas de un hospital, el despacho de energía o el mantenimiento industrial, y revisa sus permisos y su retroalimentación con OODA y la privacy-security frontier.

\end{document}
